\chapter{Premisas}
\label{ch:premises}

Un diseño es un conjunto de decisiones tomadas bajo restricciones, y la mayoría de
los malos diseños para Cuba son malos porque relajan calladamente alguna
restricción que no se relaja. Este capítulo fija el conjunto de restricciones.
Cada punto es un hallazgo de las Partes I y II y no un supuesto, cada uno se
enuncia en la forma que importa para los capítulos que siguen, y cada uno es algo
con lo que el diseño tiene que trabajar y no algo que pueda rodear.

\section{La población es la escasez vinculante}

Cuba tiene menos de diez millones de residentes y sigue perdiendo. La fecundidad
está por debajo del reemplazo desde 1978, de modo que no hay ninguna cohorte en
camino; más de una cuarta parte de la población tiene más de sesenta años; y entre
uno y dos millones de personas, desproporcionadamente en edad de trabajar y de
tener hijos, se fueron entre 2021 y 2025.

La consecuencia es la inversión del supuesto que hace la mayor parte de la
planificación del desarrollo. Cuba no es un país con exceso de mano de obra en
busca de capital que la emplee. Es un país con una fuerza laboral que se encoge y
envejece, y toda propuesta de la Parte III hay que leerla como competencia por la
misma gente escasa. Un hotel que abre saca personal de un hospital. Un sector de
trabajo remoto que crece saca ingenieros del Estado. Por eso este libro no propone
sustitución de importaciones intensiva en mano de obra, por eso el
capítulo~\ref{ch:welfare} trata la aritmética de las pensiones como su problema más
difícil, y por eso el capítulo~\ref{ch:talent} trata la inmigración y la migración
de retorno como política fiscal y no como cuestión sentimental.

Fija además el reloj. Ninguna política adoptada en 2026 cambia la población en edad
de trabajar antes de la década de 2040 salvo por migración, y la emigración se
refuerza a sí misma: cada salida elimina una razón para que la siguiente persona se
quede y añade un hogar en el exterior capaz de financiar su viaje. La primera fase
del diseño no se juzga por el crecimiento. Se juzga por si el flujo de salida se
frena.

\section{El Estado es insolvente y no tiene acceso a crédito}

Cuba está en impago de su deuda en moneda convertible desde 1986. El arreglo del
Club de París de 2015, que condonó buena parte de los atrasos acumulados y
reprogramó el resto, se rompió en cuatro años cuando Cuba dejó de pagar bajo él.
Rusia condonó el noventa por ciento de la deuda de la era soviética en 2014. Hay
proveedores con facturas impagas que se remontan a 2009, cuando el Estado congeló
las cuentas de las empresas extranjeras. Cuba no es miembro del Fondo Monetario
Internacional, el Banco Mundial ni el Banco Interamericano de Desarrollo, de modo
que no tiene acceso ni a sus préstamos ni a la asistencia técnica y la
certificación implícita que otros países usan para atraer capital privado.

De ahí se siguen dos cosas. Nada de la Parte III puede financiarse con
endeudamiento soberano en condiciones normales, porque tales condiciones no
existen; el financiamiento tiene que venir del capital propio, de la diáspora, de
fuentes concesionales y bilaterales, y sobre todo del flujo de caja. Y la secuencia
del capítulo~\ref{ch:money} tiene que poner temprano la normalización de la deuda,
no porque saldarla libere dinero ---no lo hace--- sino porque hasta que se salde
toda transacción del país se cotiza como si la contraparte incumpliera, cosa que ha
hecho.

\section{Nada funciona hasta que funcione la electricidad}

Esta es la razón de que el capítulo~\ref{ch:energy} preceda a todos los capítulos
sectoriales.

La electricidad no es un sector entre varios en Cuba; es el insumo de todos ellos.
El agua se bombea. La comida se refrigera. Los hospitales no pueden programar
salones sin ella. El turismo no puede venderse sin ella ---un hotel sin
electricidad ni agua no simplemente decepciona a un huésped, destruye la reseña que
produce al huésped siguiente---. El trabajo remoto es imposible sin ella. La
manufactura es imposible sin ella. Y una red que ha colapsado a escala nacional
cuatro veces colapsa con más facilidad cada vez.

La premisa de diseño es por tanto contundente: todo plan cuyos capítulos
sectoriales funcionarían solo si la electricidad fuera fiable es un plan que tiene
que entregar primero electricidad fiable, y sus primeros dieciocho meses pertenecen
a eso y a muy poco más.

\section{El piso de bienestar no baja}

Esto se enuncia en el prefacio y se repite aquí como premisa porque obliga a todo
lo que viene después.

El diseño conserva un Estado de bienestar universal al nivel de Europa occidental
---salud, educación, pensiones, discapacidad, sostén de ingresos--- y el resto de
la Parte III se ordena en torno a poder pagarlo. El razonamiento está en el
capítulo~\ref{ch:welfare}. Lo que corresponde aquí es el costo, admitido por
adelantado en vez de descubierto en un capítulo posterior: una estructura europea
de prestaciones exige una carga tributaria muy por encima de la norma
latinoamericana, lo que es una desventaja competitiva real frente a todos los demás
países de la región que pujan por la misma inversión. Este libro resuelve esa
disyuntiva a favor del Estado de bienestar y compensa por el lado de la inversión
con certidumbre, rapidez y Estado de derecho en vez de con tasas bajas. El
capítulo~\ref{ch:capital} está escrito sobre esa base y no debe suponer calladamente
lo contrario.

El piso está además, en 2026, más delgado que nunca y depreciándose. Conservarlo no
es cuestión de dejar algo en paz. Exige un gasto que hoy no existe, en personal que
hoy se está yendo.

\section{La diáspora es un activo con un agravio}

Algo más de dos millones de personas de origen cubano viven en el exterior, la
mayoría a tres horas de vuelo, con capital, competencias profesionales, redes de
negocio y acceso a mercados que ninguna política puede comprar de otro modo. Sus
remesas están ya entre las mayores fuentes de divisas del país. Son la mayor fuente
potencial individual de la inversión que necesita el capítulo~\ref{ch:capital}, y
son el primer mercado natural para los visitantes del capítulo~\ref{ch:tourism} y
para los clientes del capítulo~\ref{ch:talent}.

Tienen también reclamaciones y tienen memorias, y las memorias son específicas y no
difusas: casas confiscadas, el inventario levantado en la puerta, los años de
trabajo agrícola que fueron el precio de una solicitud de salida, los niños de
Peter Pan, las turbas frente a la puerta en 1980, y los parientes que no lograron
salir. Un diseño que trate a la diáspora como una chequera fracasará, y un diseño
que trate toda reclamación como válida y todo activo como devolvible fracasará de
otra manera y peor. El capítulo~\ref{ch:capital} propone un arreglo acotado y no
restitutorio, y el capítulo~\ref{ch:polity} propone un mecanismo de esclarecimiento
sin poder acusatorio, precisamente porque estas dos cosas hay que resolverlas
juntas y ninguna de las dos se resuelve con generosidad en una sola dirección.

\section{Los poderes establecidos están sentados a la mesa}

Esta es la premisa que la mayoría de las propuestas de reforma para Cuba omite, y
omitirla es lo que las vuelve poco serias.

Cualquier transición desde el arreglo actual la diseña y la ejecuta, o como mínimo
la consiente, tres grupos que hoy tienen el poder: el aparato del Partido
Comunista, las Fuerzas Armadas Revolucionarias, y el conglomerado empresarial
vinculado a los militares, GAESA, que controla una parte grande de los ingresos en
divisa del país a través del turismo, el comercio minorista y la logística ---qué
parte no se publica, que es en sí mismo el punto---.

Estos no son abstracciones que se puedan legislar hacia la desaparición. Tienen los
activos, las armas, las cuentas en divisa y la capacidad administrativa, y han
visto lo que le pasó a la nomenklatura soviética, a Ceaușescu y a Gadafi. Un diseño
que no les ofrezca nada es un diseño que no permitirán, y el registro histórico es
inequívoco en cuanto a que pueden impedirlo. Un diseño que les ofrezca todo es un
diseño que produce una oligarquía con bandera, que el capítulo~\ref{ch:integrity}
sostiene que es el modo de fallo más probable y el que resulta irreversible en la
práctica una vez ocurrido.

El libro trata por tanto a los poderes establecidos como una parte con la que hay que
negociar explícitamente: garantías de seguridad personal, pensiones y propiedad a
cambio de la entrega del control discrecional, con el trato escrito, secuenciado y
verificable. Los capítulos~\ref{ch:polity}, \ref{ch:capital} y \ref{ch:integrity}
llevan cada uno una parte de él.

\section{Estados Unidos es una variable, no una constante}

El embargo puede relajarse, endurecerse o quedarse donde está, y la política cubana
no controla cuál de las tres. La Helms-Burton lo convirtió de instrumento ejecutivo
en estatuto, de modo que su retirada exige un acto del Congreso y no una decisión
presidencial, y la política interna norteamericana de ese acto lleva treinta años
inmóvil y no da señales evidentes de moverse.

La premisa de diseño es simétrica y es la única defendible: \emph{no suponer nada}.
Nada de la Parte III puede depender de que el embargo termine, porque puede no
terminar; y nada de la Parte III puede construirse sobre el supuesto de que
persistirá, porque una normalización repentina también es posible y un país sin
plan para ella quedará desbordado por la entrada resultante ---de capital, de
visitantes, de reclamaciones--- exactamente como quedó desbordado en pequeño en
2015--16. Cada capítulo que toca la relación norteamericana enuncia por tanto qué
hace bajo cada una de las dos ramas.

\section{La raza y la región deciden quién está expuesto}

La desigualdad cubana tiene un color y un rumbo en la brújula, y a las dos las hizo
la política.

La economía de las remesas fluye abrumadoramente hacia hogares cuyos parientes se
fueron en los años sesenta, que eran abrumadoramente blancos. Los empleos en divisa
del sector turístico se han asignado en parte por apariencia. Entretanto las cinco
provincias orientales ---las más pobres, las más agrícolas, las más
afrocubanas--- tienen la peor vivienda, la peor infraestructura y los apagones más
largos, y son el origen de la mayor parte de la migración interna hacia La Habana,
que el Estado restringe por decreto desde 1997.

Toda medida liberalizadora de la Parte III favorece al hogar con capital, con un
pariente en el exterior, con una casa en un barrio turístico, con una computadora.
Eso no es un argumento contra las medidas; es un argumento para saber de antemano
quién está del otro lado de ellas, y para que los instrumentos compensatorios ---el
piso universal del capítulo~\ref{ch:welfare}, las reglas de asignación regional de
los capítulos~\ref{ch:energy} y \ref{ch:tourism}, y la estructura deliberadamente
amplia de la privatización del capítulo~\ref{ch:capital}--- estén diseñados desde el
principio y no añadidos cuando la disparidad se vuelva visible. Cuba ya corrió el
experimento de declarar resuelta la cuestión racial y luego no medirla durante
treinta años. El primer requisito del diseño en este dominio es publicar los datos.

\section{Lo que este libro no supone}

Por último, las premisas negativas, enunciadas para que un lector pueda exigirle
cuentas al argumento.

\emph{Nada de terapia de choque.} El libro no propone la liberalización simultánea
de todos los precios, ni la privatización masiva rápida, ni la retirada del Estado
social antes de que existan las instituciones que lo sustituirían. El registro
histórico de ese paquete es malo en todas partes y aquí sería peor, por la razón
demográfica de la primera sección: un país cuyos jóvenes ya se fueron no sobrevive
a una década de depresión de transición.

\emph{Nada de restitución por desalojo.} Ningún cubano es sacado de una vivienda
para devolvérsela a un propietario anterior. El capítulo~\ref{ch:capital} salda las
reclamaciones en dinero y papel, dentro de un tope.

\emph{Ningún supuesto sobre Washington.} Como arriba, en las dos direcciones.

\emph{Ningún supuesto de colapso del régimen, y ninguno de su permanencia.} El
diseño está escrito de modo que sus fases tempranas pueda ejecutarlas el gobierno
actual si lo decide, y las tardías el gobierno que le siga. Es deliberado. Un plan
que exige una ruptura política determinada antes de su primera medida es un plan
que no hace nada hasta la ruptura, y los cubanos llevan esperando la ruptura desde
1991.

\emph{Ninguna pretensión de certeza.} El capítulo~\ref{ch:sequence} cierra con una
lista de indicadores por los que puede juzgarse que el diseño está fracasando. Una
propuesta que no puede falsarse no es una propuesta.
